\section{引言}
\label{sec:intro}

再生数监测中追加病例的价值取决于所估计的量。当期传播水平、已观测区间内传播水平的平均值，以及生成连续时期的过程均值，是不同的推断目标。它们都可能服务于监测，却赋予跨期变异不同的角色。估计当期值时，传播变化属于信号；估计过程均值时，这种变化贡献抽样变异。因此，关于精度极限的陈述需要同时指明目标与观测模型，单凭病例量无法确定这一极限。

再生数概括指定条件下的传播能力，为流行态势评估提供依据\citep{fraser2007estimating, gostic2020practical}。更新方法通过代际间隔分布把发病数与既往感染压力联系起来，其不确定性取决于观测计数包含的信息\citep{cori2013new, thompson2019improved}。有关低发病率、报告噪声与及时发现传播变化的研究，已阐明逐期推断的重要限制\citep{parag2021improved, parag2022quantifying, parag2025asymmetric}。时间聚合与观测延迟又使传播过程与观测序列的表观再生数产生差异\citep{nash2023estimating, ogigittins2024simulation, eales2024differences}。这些研究促使我们区分两类不确定性：一类关于已经实现的传播轨迹，另一类关于该轨迹所抽自的时期总体。

本文在单代分枝实验中展开这一区分。原发病例独立产生负二项后代，每例继发感染独立报告，原发病例数由设计固定或外生给定。在报告概率固定且已知时，尺度化计数是当期再生数的条件无偏估计，其期望平方误差与病例量成反比。对固定的一组已观测时期取平均，若目标为这些时期的实现均值，同样可以降低计数误差。估计平稳过程均值则还需计入这些时期的抽样方差。在相同时期内增加病例不会改变这一附加项，而增加时期可以按其相依结构降低该项。

这一区分建立在经典方差结构之上。种群生态学区分个体独立结局与种群共有的环境变异\citep{engen1998demographic, lande2003stochastic}；随机流行病模型也区分种群统计贡献与环境贡献\citep{breto2009time, he2010plug}。共同条件可以包括气候、接触模式与学校校历\citep{cauchemez2008estimating, shaman2009absolute}。在抽样理论中，增加群内观测降低群内误差，增加群数则应对群间变异\citep{cochran1977sampling, hayes1999simple}。本文的贡献是为再生数监测明确对应项，并展示改变推断目标如何改变相关风险。由此得到的平衡规模是在指定模型下比较两项贡献的参照，并非普适的最低病例要求。

报告过程需要单独讨论。感染的独立漏报与同一时期所有感染共同受到的报告概率变化并不相同\citep{azmon2014estimating, stoner2019hierarchical}。即使报告概率稳定，只要其数值未知或计数归一化错误，也可能产生偏差。当报告概率独立于传播而波动时，以已知平均报告率归一化仍会给当期值估计留下额外误差；它也影响有限时期实现均值。对于过程均值，仅凭计数均值与方差一般无法分离报告波动和传播变异。这与时间序列中过程变异和观测变异难以区分的更普遍问题一致\citep{knape2008estimability, augermethe2016state}。

常规医院监测并不直接观测分枝实验中的变量。住院发生在感染之后，带有延迟，汇总了异质人群，并按日历周期记录。住院计数规模是流行过程的结果，并非外生设计变量。因此，本文考察美国国家医疗安全网络（NHSN）流感住院数据中表观周增长统计量的相对偏差平方\citep{cdc2026nhsn, frutos2025evaluation}。这些偏差描述统计量围绕经验中心的变异，并不是相对于已观测传播真值的平方误差。它们与计数倒数之间的仿射关联可以描述这种变异，却不据此识别生物学来源，也不确立观测计数范围之外的渐近极限。

实证分析区分数据能够回答的不同问题。仿射曲线与幂律曲线的比较涉及描述性形态与预测；受控二项稀释考察对既有观测轨迹追加抽样所引起的响应；州截距拟合考察州内关联，替代中心则评估共同时间模式带来的敏感性。模拟在给定代际间隔、延迟与相依结构的条件下，将这些摘要与指定的传播及观测机制联系起来。上述任何操作都不能单独识别传播环境方差。其作用、不确定性计算与执行状态在方法和结果中说明（第~\ref{sec:methods}~节与第~\ref{sec:results}~节）。

因此，本文先明确推断目标，再解释病例量的作用。第~\ref{sec:background}~节给出精确恒等式与矩不可识别结果，表~\ref{tab:targets} 汇总各目标。随后，实证分析在观测统计量尺度上评价关联及其敏感性。证明以及预测、收缩估计与设计的符号化扩展置于附录。各部分共同说明：追加病例在何种条件下改善所选估计，以及哪些剩余不确定性需要病例计数以外的信息。
